\documentclass[11pt]{article}
\usepackage[T1]{fontenc}
\usepackage{amsmath,amssymb,amsthm}
\usepackage{hyperref}
\usepackage{microtype}
\usepackage{booktabs}
\usepackage[italian]{babel}
\emergencystretch=3em

% Rimuove il quadratino di fine dimostrazione
\renewcommand{\qedsymbol}{}

\title{Prova di Continuità: un modello temporale\\
per la propagazione dell'autorità nei\\
sistemi distribuiti e negli agenti di IA}
\author{Nicola Gallo}
\date{8 luglio 2026}

\newtheorem{definition}{Definizione}
\newtheorem{lemma}{Lemma}
\newtheorem{theorem}{Teorema}
\newtheorem{corollary}{Corollario}

\begin{document}
\maketitle

\begin{abstract}
I modelli di autorizzazione basati sulla Prova di Possesso fanno discendere
l'autorità dal possesso di artefatti quali token, credenziali o capability. In
questo lavoro sosteniamo che il possesso non sia sufficiente per le catene di
esecuzione discrete, tanto quando attraversano più servizi quanto quando si
articolano in passi distinti all'interno della stessa macchina: esso, infatti,
non garantisce che si conservi la relazione causale fra l'origine di una
richiesta e l'autorità esercitata nei passi successivi. Introduciamo la Prova di
Continuità, una disciplina minimale per la propagazione dell'autorità nel
modello Provenance Identity Continuity (PIC). In tale disciplina, ogni passo di
esecuzione deve essere legato causalmente al precedente e può propagare soltanto
un sottoinsieme non espansivo dell'autorità ricevuta dall'origine. Il modello
introduce inoltre la Prova di Relazione, una primitiva causale relativa a un
singolo passaggio, la cui composizione transitiva costituisce la Prova di
Continuità; entrambe integrano la Prova di Possesso, senza sostituirla. Nel
modello proposto, la condizione del delegato confuso non può verificarsi come
comportamento valido: ogni privilegio esercitato in un passo successivo deve
essere già presente nel contesto di autorità dell'origine. Il risultato è
direttamente pertinente ai sistemi distribuiti e agli agenti di IA, nei quali
gli esecutori invocano strumenti e servizi a valle pur disponendo di più fonti
di autorità, cosicché il medesimo disallineamento fra autorità e causalità si
ripresenta oltre i confini dei servizi. Con la Prova di Continuità queste fonti
possono essere trasportate insieme, ma non vengono mai fuse in un'unica autorità
combinata, poiché ciascun passo è autorizzato esclusivamente rispetto al
contesto di autorità della linea causale che lo ha determinato.

Il lavoro riguarda la propagazione dell'autorizzazione, non l'autenticazione:
meccanismi di identità e autenticazione quali OIDC, credenziali verificabili,
wallet e identità dei carichi di lavoro restano strumenti complementari per
stabilire l'origine; la Prova di Continuità disciplina invece il modo in cui
l'autorità si propaga una volta stabilita tale origine.
\end{abstract}

\section{Modello}

Formalizziamo il modello PIC (Provenance Identity Continuity) come segue.

Siano $P$ un insieme finito di soggetti, $O$ un insieme di operazioni e $R$ un
insieme di risorse. Definiamo privilegio ogni coppia $(o, r) \in O \times R$.

A ogni soggetto $p$ è associato un insieme di privilegi:
\[
Priv(p) \subseteq O \times R.
\]

Qui un soggetto è un'entità astratta dotata di permessi: può trattarsi di
un'identità umana stabilita mediante un fornitore di identità, un wallet, un
flusso sul modello di OIDC/OIDC4VC~\cite{openidconnect2014,openid4vc2024} o un
identificatore decentralizzato~\cite{didcore2022}; di un'identità di carico di
lavoro o di macchina, come quelle impiegate negli ambienti di tipo
WIMSE/SPIFFE~\cite{wimse2024,spiffe2017}; di un ruolo; di un account di servizio;
oppure di un'altra entità autenticata e dotata di permessi. Ai fini del modello
contano due proprietà: che all'entità sia associato un insieme di privilegi
$Priv(p)$ e che essa possa esprimere un intento entro tali privilegi. Quando
un'entità $p$ esprime un simile intento, seleziona un sottoinsieme di $Priv(p)$;
questo sottoinsieme diventa il contesto di autorità iniziale
$C_0 \subseteq Priv(p)$ di una nuova linea causale di esecuzione. I privilegi in
$Priv(p)$ possono essere espressi al livello delle operazioni che il soggetto è
autorizzato a \emph{causare}, non soltanto di quelle che compie direttamente;
un passaggio a valle può realizzare tale autorità impiegando un vocabolario
locale di operazioni differente, mediante la traduzione di policy $\mathcal{T}$
della Sezione~\ref{sec:poc}.

Supponiamo, per esempio, che Bob disponga dei privilegi
$(\text{read},\text{Book})$ e $(\text{write},\text{Book})$. Quando esprime un
intento, Bob ne seleziona un sottoinsieme --- poniamo il solo
$(\text{write},\text{Book})$ --- che diventa il contesto di autorità iniziale
$C_0$ di una nuova linea causale. Due esecuzioni possono coinvolgere entrambe
Bob e il medesimo privilegio selezionato senza per questo costituire la stessa
occorrenza: ciascuna trasporta l'autorità propagata lungo la propria linea
causale. Una può essere la continuazione valida di un intento espresso da Bob
entro una data linea; l'altra può essere un'esecuzione indipendente che esercita
lo stesso privilegio apparente al di fuori della linea che lo ha conferito. Le
coordinate dell'identità e del privilegio possono coincidere; è la coordinata
della linea causale a distinguere le due esecuzioni.

Una richiesta è un messaggio $req$, eventualmente dotato di un carico utile,
trasmesso fra esecutori, cioè le entità che compiono i passi di esecuzione.

L'esecuzione è modellata come una catena causale di passi:
\[
\pi = \langle (s_0, C_0), (s_1, C_1), \dots, (s_n, C_n) \rangle
\]
ove $s_0$ è il passo iniziale e $C_i \subseteq O \times R$ è il contesto di
autorità trasportato al passo $s_i$. Indichiamo con $s_i$ il passo di esecuzione
formale al passaggio $i$; sul piano operativo, un passaggio consiste
nell'esecuzione di uno di tali passi da parte di un esecutore che riceve,
elabora ed eventualmente propaga una richiesta. Ogni passaggio può essere
compiuto da un servizio, un carico di lavoro, una funzione, uno strumento o un
agente; il modello richiede la continuità dell'autorità, non la propagazione
dell'identità.

\begin{definition}[Modello PIC]
Un sistema garantisce la Provenance Identity Continuity se, per ogni catena di
esecuzione $\pi$, ciascuna transizione conserva la continuità causale e la
restrizione monotona dell'autorità:
\[
C_{i+1} \subseteq C_i
\]
per ogni $i$, e se ogni privilegio esercitato al passaggio $n$ è collegato
causalmente all'origine mediante una catena verificabile.
\end{definition}

Il legame causale relativo a un singolo passaggio e attestato a ogni transizione
è una \emph{Prova di Relazione} ($PoR$): un'evidenza che certifica come un passo
sia la continuazione valida del suo predecessore immediato entro la stessa linea
causale. Il modello PIC combina dunque due condizioni --- il legame causale,
attestato passaggio per passaggio da $PoR$, e la restrizione monotona
$C_{i+1}\subseteq C_i$ ---, qui presentate informalmente e definite in modo
formale nella Sezione~\ref{sec:poc}, dove la continuità è costruita come
composizione transitiva di tali relazioni. Il teorema seguente enuncia la
conseguenza della seconda condizione sull'autorità.

\begin{theorem}[Sicurezza PIC]
Se vale il modello PIC, nessun passo di esecuzione può esercitare un privilegio
assente all'origine.
\end{theorem}

\begin{proof}
Per costruzione, $C_{i+1} \subseteq C_i$ per ogni $i$, dunque
$C_n \subseteq C_0$. Ne segue che ogni $(o,r) \in C_n$ appartiene anche a
$C_0$. Non è quindi possibile acquisire alcun privilegio ulteriore rispetto al
contesto di autorità iniziale.
\end{proof}

\section{Modello di minaccia e ambito}

Consideriamo catene di esecuzione nelle quali una richiesta iniziale determina
uno o più passi di esecuzione successivi. Gli esecutori possono contenere
errori, trovarsi in uno stato di confusione oppure subire un'influenza avversaria
attraverso i propri input o il proprio contesto di esecuzione. Un esecutore può
disporre di più fonti di autorità --- privilegi ambientali, token delegati,
credenziali al portatore o capability --- e può tentare di usarne una estranea
alla richiesta che sta elaborando.

L'avversario può influenzare le richieste, i carichi utili e i parametri
trasmessi lungo la catena, inducendo un esecutore a scegliere la fonte di
autorità errata. Non può però falsificare l'evidenza di relazione richiesta da
$PoR$, né indurre un verificatore di continuità valido ad accettare una
transizione che viola la monotonia. Così come i modelli basati su capability e
token assumono che le proprie credenziali non siano falsificabili, assumiamo non
falsificabili l'evidenza di relazione e il meccanismo di transizione
dell'autorità; la compromissione del validatore, della radice crittografica di
fiducia o dell'inizializzazione dell'autorità esula dall'ambito del lavoro.

\section{Delegato confuso}

\begin{definition}[Delegato confuso]\label{def:confused}
Si verifica un caso di delegato confuso quando esistono due soggetti $U$
(utente) e $D$ (delegato) e un privilegio $(o,r) \in O \times R$ tali che:
\begin{enumerate}
    \item $(o,r) \notin Priv(U)$,
    \item $(o,r) \in Priv(D)$,
    \item $U$ invia una richiesta $req$ a $D$,
    \item in conseguenza di $req$, $D$ esegue $(o,r)$.
\end{enumerate}
\end{definition}

La definizione non dipende dai dettagli implementativi, ma soltanto dal
disallineamento fra autorità e causalità.

\paragraph{Esempio classico.}
Un compilatore FORTRAN \texttt{FORT}, installato nella directory privilegiata
\texttt{SYSX}, dispone dell'autorità ambientale necessaria a scrivere le
statistiche in \texttt{(SYSX)STAT}. Un utente invoca il compilatore e indica
\texttt{(SYSX)BILL} come nome del file di output per il debug. Il compilatore
apre in scrittura il file di destinazione avvalendosi della propria autorità su
\texttt{SYSX}, sovrascrivendo così \texttt{(SYSX)BILL}. L'utente non ha mai
posseduto tale privilegio, ma il delegato lo ha esercitato in conseguenza della
richiesta dell'utente~\cite{hardy1988confused}.

\section{Prova di Possesso (PoP)}

Un token $t$ rappresenta un insieme di privilegi:
\[
Priv(t) \subseteq O \times R.
\]

\paragraph{Semantica di PoP.}
Il possesso implica la possibilità d'uso: se un soggetto possiede $t$, può
esercitare ogni $(o,r)\in Priv(t)$.

I sistemi PoP vincolano l'autorità al possesso di un artefatto, ma non
richiedono, di per sé, che l'autorità esercitata sia una continuazione causale
del contesto di autorità della richiesta.

\subsection{Token vincolati al mittente e al possessore}

I token vincolati al mittente o al possessore --- per esempio
DPoP~\cite{fett2023dpop} e Token Binding~\cite{popov2018token} --- dimostrano
\emph{chi} possa presentare o usare un artefatto, legandolo a una chiave o a un
canale di trasporto. Non dimostrano però, da soli, che l'autorità esercitata
tramite tale artefatto continui l'attuale linea causale di esecuzione. Limitare
chi possa presentare un token pone un vincolo sul trasporto, non sulla
causalità: in un dato passo, un possessore legittimo può comunque esercitare
un'autorità in suo possesso ma esterna alla linea causale della richiesta che
sta elaborando. Il vincolo risponde alla domanda sull'autenticità del
possessore; resta invece aperta la questione causale, ossia se questo uso
prosegua il contesto di autorità che lo ha determinato.

\subsection{Commistione di autorità}

Supponiamo che:
\begin{align}
&(\text{write},r) \notin Priv(U) \tag{H1}\\
&(\text{write},r) \in Priv(D) \tag{H2}
\end{align}

Qui $U$ è l'utente e $D$ un delegato.

In un modello basato sul possesso, un delegato può disporre di più fonti di
autorità: i propri privilegi ambientali, token delegati, credenziali al
portatore o altri artefatti. Il modello autorizza un'operazione verificando se
una delle fonti di autorità possedute concede il privilegio richiesto.

\begin{definition}[Funzione di selezione]
Sia $\{A_j\}$ la famiglia delle fonti di autorità detenute da un esecutore ---
privilegi ambientali, token delegati e artefatti vincolati al possessore o al
mittente. Una \emph{funzione di selezione} $\sigma$ sceglie, per un'operazione
richiesta $(o,r)$, una fonte $A_j$ tale che $(o,r) \in A_j$, indipendentemente
dalla linea causale $\ell$ della richiesta --- ossia la linea causale della
richiesta in corso di elaborazione, formalizzata nella
Sezione~\ref{sec:poc}.
\end{definition}

\begin{theorem}[PoP ammette la commistione di autorità]
In un modello basato sul possesso, se un esecutore può selezionare l'autorità da
un qualsiasi artefatto o privilegio ambientale di cui dispone,
indipendentemente dal contesto di autorità della richiesta, allora il modello
ammette la condizione del delegato confuso.
\end{theorem}

\begin{proof}
Consideriamo una richiesta che ha origine in un contesto di autorità privo di
$(o,r)$. Supponiamo che l'esecutore disponga di un altro contesto di autorità,
artefatto o privilegio ambientale che invece contiene $(o,r)$.

In un modello PoP, l'autorizzazione dipende dal possesso di un artefatto o di un
privilegio che concede $(o,r)$, non dal fatto che tale autorità continui il
contesto di autorità della richiesta. Il modello consente dunque un'esecuzione
nella quale la richiesta determina causalmente $(o,r)$ benché $(o,r)$ sia
assente dal contesto di autorità della richiesta. Questa è precisamente la
condizione del delegato confuso della Definizione~\ref{def:confused}.
\end{proof}

Il ragionamento vale per ogni modello nel quale l'autorizzazione discende dal
possesso indipendentemente dalla linea causale, compresi i JWT, le credenziali
al portatore e le capability trattate come artefatti posseduti quando possono
essere applicate al di fuori dell'esecuzione che le ha richieste. Il vincolo al
possessore e quello al mittente limitano quale soggetto possa presentare $A_j$,
non stabiliscono se la scelta operata da $\sigma$ appartenga a $\ell$; pertanto
non eliminano tale ammissibilità. La lacuna non risiede nel possesso in sé, ma
nell'assenza di un vincolo temporale fra l'autorità impiegata e l'esecuzione che
ne ha causato l'uso.

Si tratta di condizioni al livello dell'esecuzione: l'autorità appartenente a
una linea causale viene introdotta in un'altra, dando luogo al delegato confuso.
La Prova di Continuità elimina questa ammissibilità al livello del modello: in
un'esecuzione PIC valida la condizione del delegato confuso non può verificarsi
(Sezione~\ref{sec:cdimposs}), poiché l'autorità è definita esclusivamente come
continuazione della linea causale che l'ha determinata.

\section{Prova di Continuità (PoC / PIC)}
\label{sec:poc}

L'esecuzione è una catena causale:
\[
s_0 \rightarrow s_1 \rightarrow \dots \rightarrow s_n.
\]

Indichiamo con $C_i \subseteq O \times R$ il contesto di autorità trasportato dal
passo di esecuzione $s_i$.

\begin{definition}[Prova di Relazione]
$PoR(s_i, s_{i+1})$ vale se e solo se $s_{i+1}$ è una continuazione causale
valida di $s_i$ entro la medesima linea causale di esecuzione. $PoR$ è una
proprietà \emph{relazionale}, riferita a un singolo passaggio: lega ciascun
passo al suo predecessore immediato. Il collegamento retrospettivo con l'origine
non è contenuto in una singola $PoR$, ma si ottiene componendo relazioni
consecutive, come precisa la definizione di $PoC$ riportata di seguito.
\end{definition}

\begin{definition}[Transizione valida]
Una transizione da $s_i$ a $s_{i+1}$ è valida se e solo se è al tempo stesso
legata causalmente e non espansiva:
\[
Valid(s_i,C_i,s_{i+1},C_{i+1})
\Longleftrightarrow
PoR(s_i,s_{i+1}) \land C_{i+1}\subseteq C_i.
\]
\end{definition}

La prima condizione dimostra la relazione causale; la seconda dimostra la
restrizione monotona dell'autorità. Ogni passaggio trasferisce quindi soltanto:
\[
C_{i+1} \subseteq C_i. \tag{C1}
\]

\begin{definition}[Prova di Continuità]
Per una catena $\pi = \langle s_0, \dots, s_n \rangle$, $PoC(\pi)$ vale se e
solo se ogni transizione fra passi adiacenti è valida:
\[
PoC(\pi) \Longleftrightarrow
\bigwedge_{i=0}^{n-1} \big( PoR(s_i, s_{i+1}) \land C_{i+1} \subseteq C_i \big).
\]
\end{definition}

La continuità è dunque la composizione transitiva delle relazioni: $PoR$
dimostra un singolo legame, $PoC$ l'intera linea causale. La relazione è locale;
la continuità è globale.

\paragraph{Propagazione monotona lungo una catena causale.}
Due elementi si propagano congiuntamente. In primo luogo, l'autorità si attenua
in modo monotono: scriviamo $C_{i+1} \preceq C_i$ per indicare
l'\emph{attenuazione} dell'autorità --- ogni autorità presente a un passaggio è
una restrizione di quella presente al passaggio precedente, mai un'espansione.
Quando il contesto è un insieme di privilegi atomici, la relazione si
specializza in $C_{i+1} \subseteq C_i$, che è la forma usata nei nostri esempi.
Sebbene il nucleo del modello rappresenti i contesti di autorità come insiemi
di privilegi atomici, usiamo $\preceq$ per riferirci all'ordine generale di
attenuazione; nel caso insiemistico adottato in tutti gli esempi formali, tale
ordine coincide con $\subseteq$. In secondo luogo, $PoR$ lega i passaggi in
un'unica linea causale: la continuità è proprietà non di un passaggio isolato,
bensì della \emph{catena}. Formalmente, la linea è una catena causale logica;
un'implementazione non deve necessariamente trasportarla per intero: può
rappresentarla mediante commitment crittografici, riferimenti, hash, firme o
prove sintetiche che consentano a ogni passaggio di dimostrare la propria
posizione nella linea senza trasportarne l'intera storia.

\paragraph{Inserimento dell'esecuzione nel modello di autorità.}
La Prova di Continuità è definita proprio per questo scenario: non è necessario
che il passo successore sia noto, scelto, istanziato o predisposto quando il
predecessore termina il proprio passo. L'esecuzione è una sequenza temporale di
passi causali, non di posizioni fissate dalla topologia: l'esecutore a monte
agisce a un istante $t$, mentre quello a valle può venire a esistenza soltanto a
un istante successivo $t' > t$. Non è quindi sempre possibile vincolare in
anticipo l'autorità a uno specifico possessore, a una chiave o al canale di un
successore ancora ignoto.

Ciò che deve conservarsi oltre tale discontinuità non è l'identità del
successore, né una policy a esso associata, bensì il fatto che la sua autorità
continui il contesto che lo ha determinato. Una credenziale, un'attestazione
dell'identità del carico di lavoro, una credenziale di account di servizio o un
token vincolato possono autenticare il possessore oppure legare un artefatto a
una chiave, a un carico di lavoro o a un canale. Tuttavia, se la decisione di
autorizzazione non considera anche evidenze sensibili alla linea causale, i
medesimi attributi del possessore possono presentarsi in esecuzioni distinte:
essi stabiliscono chi agisce, non quale catena causale abbia prodotto la
richiesta. Derivare nuovamente l'autorità in $i+1$ da un'asserzione nuova sul
successore, anziché dalla sua relazione con il passo che lo ha determinato,
reintroduce la lacuna di individuazione della
Sezione~\ref{sec:impossibility} --- la condizione strutturale alla base del
delegato confuso (Definizione~\ref{def:confused}). La continuità evita il
problema perché non può essere autoasserita: essendo definita soltanto in
relazione al predecessore, può essere \emph{dimostrata} come continuazione, ma
non asserita isolatamente in $i+1$.

Un predecessore può invece emettere una continuazione che sarà poi consumata da
un successore:
\[
(s_i,C_i)
\;\xrightarrow{\;\mathrm{emit}\;}
\text{unknown next hop}
\;\xrightarrow{\;PoR,\; C_{i+1}\subseteq C_i\;}
(s_{i+1},C_{i+1}).
\]
Oltre tale discontinuità persiste una continuità valida, non il possesso da
parte di un titolare preesistente: persistono una relazione causale con il
predecessore e un contesto di autorità non espansivo. La Prova di Relazione
attesta tale proprietà dell'esecuzione, anziché del possessore; composta
transitivamente nel rispetto dell'attenuazione monotona
($C_{i+1} \preceq C_i$), essa costituisce la Prova di Continuità, cioè
l'invariante di propagazione del modello PIC.

\paragraph{Spazi di operazioni eterogenei.}
La notazione presuppone un universo comune dei privilegi $O \times R$. Nei
sistemi eterogenei, passaggi diversi impiegano vocabolari differenti; sia
$O_i \times R_i$ l'universo al passaggio $i$. Una transizione è allora regolata
da una traduzione di policy
\[
\mathcal{T}_{i \to i+1} : \mathcal{P}(O_i \times R_i) \to
\mathcal{P}(O_{i+1} \times R_{i+1}),
\]
e la monotonia diventa
$C_{i+1} \subseteq \mathcal{T}_{i \to i+1}(C_i)$: l'autorità del successore non
supera quella consentita dal predecessore secondo la traduzione locale.
Consideriamo $\mathcal{T}$ un dato di policy fissato; la sintesi e la verifica di
tali traduzioni esulano dall'ambito del lavoro. L'eterogeneità è qui semantica,
non meramente amministrativa: passaggi diversi possono descrivere l'autorità
mediante vocabolari di operazioni e risorse differenti persino entro un unico
perimetro di fiducia. Nei sistemi eterogenei, la sicurezza è pertanto relativa
alla correttezza della traduzione di policy scelta
$\mathcal{T}_{i \to i+1}$; una traduzione eccessivamente permissiva costituisce
un errore di policy, non una violazione dell'invariante di continuità in sé.

\begin{lemma}[La continuità implica il vincolo all'origine]
Se vale $PoC(\pi)$, allora $s_n$ è collegato causalmente a $s_0$ mediante una
catena relazionale verificabile e $C_n \subseteq C_0$.
\end{lemma}

\begin{proof}
Ogni $PoR(s_i,s_{i+1})$ collega due passi consecutivi; la loro composizione
collega $s_n$ a $s_0$. La non espansività a ogni passaggio dà
$C_n \subseteq \dots \subseteq C_0$.
\end{proof}

Precisiamo esplicitamente l'ambito: il lavoro modella l'invariante di
continuità; la costruzione concreta (per esempio crittografica) di $PoR$ esula
da tale ambito ed è realizzata da un'architettura di applicazione complementare.
Così come i modelli basati su capability e token assumono non falsificabili le
proprie credenziali, assumiamo non falsificabile la linea causale.

\subsection{Regola di autorizzazione}

La regola seguente viene valutata soltanto per una catena di continuità valida,
ossia una catena nella quale ogni transizione fra passi adiacenti soddisfa sia
$PoR(s_i,s_{i+1})$ sia $C_{i+1}\subseteq C_i$.

Nell'universo comune dei privilegi $O \times R$, $(o,r)$ è autorizzato al
passaggio finale se e solo se:
\[
(o,r) \in \bigcap_{i=0}^{n} C_i.
\]

Poiché la catena è monotona decrescente:
\[
C_n \subseteq C_{n-1} \subseteq \dots \subseteq C_0,
\]
si ha:
\[
\bigcap_{i=0}^{n} C_i = C_n.
\]

Nessun passaggio può quindi acquisire un'autorità che non fosse già presente
all'origine, e ogni autorità abbandonata a un certo passaggio $j$ è perduta in
modo irreversibile: per monotonia, non può ricomparire in alcun passaggio
successivo.

\subsection{Il cambiamento ontologico: possesso e continuità}

I modelli basati sul possesso e sulle
capability~\cite{dennis1966programming,miller2003capability} rispondono a una
domanda \emph{spaziale}: un attore detiene autorità su un oggetto? L'autorità è
un punto di $O \times R$. Il modello è atemporale: non possiede alcun asse sul
quale registrare l'esecuzione che ha condotto a un'azione.

La Prova di Relazione fornisce l'asse mancante, legando un passo di esecuzione
al suo predecessore causale. Su di essa si costruisce la Prova di Continuità ---
la continuità è la composizione transitiva delle relazioni ---, che combina due
condizioni: la restrizione monotona dei privilegi
($C_{i+1} \subseteq C_i$) e, mediante $PoR$, la continuità causale
dell'esecuzione che li trasporta.

L'autorità non appartiene dunque più al solo $O \times R$, bensì a
$O \times R \times \mathcal{L}$, dove $\mathcal{L}$ è la linea causale
dell'esecuzione. In concreto, $\mathcal{L}$ è l'insieme delle linee causali
lineari finite: una linea è una sequenza
$\ell = \langle s_0 \prec s_1 \prec \dots \prec s_n \rangle$, nella quale
$\prec$ denota la precedenza causale, attestata fra passi adiacenti da $PoR$;
ogni evento è così un'occorrenza di un privilegio entro una particolare linea
causale. Due azioni identiche quanto al possesso --- stesso soggetto, stesso
token, stesso privilegio --- possono differire quanto alla continuità, perché
discendono da cause diverse. Il medesimo $(o,r)$, esercitato come continuazione
valida di una richiesta autorizzante, oppure da un esecutore che attinge a
un'autorità estranea di cui casualmente dispone, è indistinguibile dal punto di
vista del possesso; si tratta tuttavia di due esecuzioni distinte: la prima
continua una linea causale autorizzante, la seconda non ne continua alcuna.
Il possesso è la proiezione su $O \times R$ che dimentica l'asse della linea
causale. Esso è la componente spaziale dell'autorità, mentre la linea causale ne
è la componente temporale; un modello basato sul solo possesso è pertanto una
proiezione dell'autorità che dimentica il tempo. Ciò non rende errati i modelli
basati sul possesso, ma li rende incompleti nel caso multipassaggio, nel quale la
continuazione dell'autorità diventa una proprietà temporale. L'evento rilevante
non è dunque il solo $(p,o,r)$, ossia identità più privilegio, bensì
un'occorrenza di un privilegio entro una linea causale. Quando il medesimo
soggetto e il medesimo privilegio ricorrono in due linee diverse, la coordinata
dell'identità non li distingue; quella della linea causale sì.

Nel linguaggio delle proiezioni, questa è la \emph{lettura temporale} del
delegato confuso della Definizione~\ref{def:confused}: un'azione valida nello
spazio ma non nel tempo --- il delegato dispone di $(o,r)$, ma
$(o,r) \notin C_0$ per la richiesta che l'ha causata ---, cioè una collisione
nella proiezione che dimentica $\mathcal{L}$: due esecuzioni distinte che il
possesso non è in grado di distinguere. La Sezione~\ref{sec:openpassthrough}
fornisce la caratterizzazione operativa (il \emph{tramite aperto}), affinando la
Definizione~\ref{def:confused} mediante la designazione.

\paragraph{Relazione con le capability.}
Le capability risolvono la questione spaziale e vanno oltre: fondendo
designazione e autorità nell'atto dell'invocazione, una capability lega
l'autorità a un singolo passo causale --- nei nostri termini, una
\emph{continuità a passaggio singolo}, cioè una Prova di Relazione ---, cosicché
il delegato confuso classico è escluso, non semplicemente mitigato. Ciò che manca
a una capability trattata unicamente come artefatto posseduto è l'estensione di
questo vincolo lungo una catena: se applicata al di fuori della linea causale che
l'ha richiesta, essa ricade nel possesso invariante rispetto alla linea.

Prova di Possesso e Prova di Continuità sono due primitive di propagazione
dell'autorità. L'invocazione di una capability è un'istanza a singolo passaggio
della primitiva di continuità, mentre una capability detenuta e presentata come
artefatto si fonda sul possesso; un sistema di capability che trasportasse la
continuità lungo l'intera catena sarebbe, in questi termini, conforme a PIC. PIC
è dunque un raffinamento temporale della propagazione basata sulle capability,
non un suo sostituto: introduce una dimensione che il solo possesso non può
codificare. Entrambi eliminano per costruzione il delegato confuso, anziché
affidarsi alla vigilanza in fase di esecuzione, ma operano su estensioni diverse.
Nel modello a capability di oggetto, la fusione di designazione e autorità al
momento dell'invocazione elimina strutturalmente il classico disallineamento a
singolo passaggio: un delegato non può applicare autorità a una destinazione per
la quale non l'ha ricevuta. Il disallineamento riappare soltanto quando si esce
dal modello a capability --- in presenza di autorità ambientale, oppure quando
una capability è usata come artefatto posseduto al di fuori dell'occorrenza che
l'ha trasportata: il caso, analizzato sopra, della capability come possesso, nel
quale l'autorità torna a dipendere dal possesso. PIC estende la stessa garanzia
per costruzione a una linea causale multipassaggio: il delegato confuso non è
uno stato rappresentabile in alcun passaggio di un'esecuzione PIC valida, perché
un privilegio assente dal contesto di autorità iniziale non costituisce affatto
uno stato successivo valido.

\subsection{Proiezione: il possesso invariante rispetto alla linea non può
distinguere le occorrenze}
\label{sec:impossibility}

Precisiamo la proiezione. Un \emph{evento} è un'occorrenza che esercita un
privilegio entro una linea causale,
\[
e = (o, r, \ell) \in E := O \times R \times \mathcal{L},
\]
e $p : E \to O \times R$, $p(o,r,\ell) = (o,r)$, è la proiezione che dimentica
la linea causale.

\begin{definition}[Policy basata sul possesso]
Una policy di autorizzazione $A : E \to \{0,1\}$ è \emph{basata sul possesso} se
e solo se è invariante rispetto alla linea causale: essa si fattorizza
attraverso $p$, cioè esiste $\bar{A} : O \times R \to \{0,1\}$ tale che
$A = \bar{A} \circ p$. Equivalentemente,
$A(o,r,\ell) = A(o,r,\ell')$ per ogni $\ell, \ell' \in \mathcal{L}$: il
verdetto non dipende dalla linea che ha prodotto l'evento.
\end{definition}

\begin{theorem}[Le policy invarianti rispetto alla linea non distinguono le
occorrenze]
Siano $e = (o,r,\ell)$ ed $e' = (o,r,\ell')$ due eventi con il medesimo
privilegio ma linee causali distinte, $p(e) = p(e')$ e $\ell \neq \ell'$. Allora
ogni policy basata sul possesso soddisfa $A(e) = A(e')$.
\end{theorem}

\begin{proof}
$A = \bar{A} \circ p$ e $p(e) = p(e')$, pertanto
$A(e) = \bar{A}(p(e)) = \bar{A}(p(e')) = A(e')$.
\end{proof}

Il delegato confuso è un caso particolare di coppia vietata: due eventi che
condividono $(o,r)$ ma differiscono per linea causale, uno dei quali deve essere
autorizzato e l'altro negato. Una linea ha origine in un contesto che contiene
$(o,r)$; l'altra esercita $(o,r)$ senza aver mai continuato un contesto che lo
contenesse. Per il teorema, \emph{nessuna} policy basata sul possesso e
invariante rispetto alla linea può distinguerli.

\begin{theorem}[Ogni soluzione reintroduce la continuità]\label{thm:reintro}
Una policy $A$ autorizza $e$ e nega $e'$, con $p(e) = p(e')$ e
$\ell \neq \ell'$, soltanto se $A$ non è invariante rispetto alla linea; in
altri termini, $A$ deve leggere una funzione $g : \mathcal{L} \to X$ della linea
causale tale che $g(\ell) \neq g(\ell')$.
\end{theorem}

\begin{proof}
Se $A$ fosse invariante rispetto alla linea, si fattorizzerebbe attraverso $p$
e, per il teorema precedente, darebbe $A(e) = A(e')$, in contraddizione con
$A(e) \neq A(e')$. Dunque $A$ dipende dalla coordinata della linea causale:
esiste una funzione $g(\ell)$ che distingue $e$ da $e'$.
\end{proof}

Una tale $g$ è esattamente un meccanismo di continuità, ossia un'estrazione
dalla linea causale che funge da prova di relazione. Le mitigazioni impiegate
nei sistemi basati sul possesso (nonce, chiavi di idempotenza, marcature
temporali, vincoli al destinatario e al contesto) sono istanze concrete di $g$:
risolvono il delegato confuso soltanto reintroducendo nella decisione
l'informazione sulla linea causale, cioè la continuità. È irrilevante che la
linea sia trasportata insieme alla richiesta o incorporata nell'artefatto: un
token che codifica un nonce o un contesto causale rende $A$ dipendente da $\ell$
e dunque, per definizione, sensibile alla continuità, non basato sul possesso.
La distinzione non riguarda il luogo in cui risiede l'informazione, bensì il
fatto che la decisione la legga o meno. Il limite non è un difetto ingegneristico
dei sistemi basati sul possesso, ma un confine modellistico dell'autorizzazione
invariante rispetto alla linea: per il primo teorema, essa è strutturalmente
incapace di distinguere le occorrenze, e ogni rimedio efficace è sensibile alla
linea --- dunque alla continuità --- anziché invariante rispetto a essa.

\paragraph{Perché il modello del possesso è rappresentativo.}
Nell'astrazione qui considerata, un token al portatore usato senza un vincolo
esplicito al contesto o alla linea causale è invariante rispetto a quest'ultima:
la decisione legge l'artefatto e il privilegio richiesto, non la linea causale
del suo uso. Un JWT, per esempio, asserisce che ``il possessore può leggere'',
non quale catena causale abbia condotto a tale lettura. Questi sistemi spesso
funzionano non perché codifichino esplicitamente la continuità, ma perché un
perimetro ne fornisce implicitamente una parte: TLS mutuo, confini di rete e
delegazione desunta dal trasporto suppliscono spesso alla linea causale mancante.
Quando il perimetro viene meno --- nelle installazioni zero-trust, nelle catene
fra più organizzazioni sulla rete pubblica e negli agenti autonomi di IA che
invocano strumenti e servizi ---, l'invarianza rispetto alla linea si manifesta
come vulnerabilità al secondo passaggio. La definizione di policy basata sul
possesso data sopra descrive dunque sistemi effettivamente impiegati, non una
loro caricatura. Un perimetro di fiducia può celare la lacuna temporale facendo
apparire autorità e causalità come co-localizzate, ma il modello non presuppone
l'attraversamento di un simile perimetro: il disallineamento può sorgere ogni
volta che l'autorità viene scelta indipendentemente dall'esecuzione che ne ha
causato l'uso.

\subsection{Un teorema di compromesso per la propagazione dell'autorità}

La proiezione della Sezione~\ref{sec:impossibility} costituisce il fondamento
tecnico di un compromesso fra tre proprietà che si può richiedere a un sistema
di propagazione dell'autorità di soddisfare simultaneamente.

\begin{definition}[Autorizzazione invariante rispetto alla linea]
L'autorizzazione di un evento $e = (o,r,\ell)$ dipende soltanto dall'autorità
posseduta per $(o,r)$, non dalla linea causale $\ell$ che ha determinato
l'evento.
\end{definition}

\begin{definition}[Delegazione che consente la commistione di autorità]
Un esecutore può disporre sia di un contesto di autorità $C_0$ derivato dalla
richiesta, sia di una fonte di autorità indipendente $C_D$, con un $(o,r)$ tale
che
$(o,r) \notin C_0$ e $(o,r) \in C_D$.
\end{definition}

\begin{definition}[Sicurezza rispetto al delegato confuso]
Nessuna richiesta il cui contesto iniziale $C_0$ non contenga $(o,r)$ può
indurre un esecutore successivo a esercitare $(o,r)$ mediante un'autorità non
derivata dalla richiesta stessa.
\end{definition}

\begin{theorem}[Compromesso fra possesso, delegazione e sicurezza]
Nessun sistema di propagazione dell'autorità può soddisfare simultaneamente
l'autorizzazione invariante rispetto alla linea, la delegazione che consente la
commistione di autorità e la sicurezza rispetto al delegato confuso.
\end{theorem}

\begin{proof}
Supponiamo che valgano tutte e tre. Per la delegazione che consente la
commistione di autorità, esistono un esecutore $D$, un contesto $C_0$ derivato
dalla richiesta, una fonte indipendente $C_D$ e un $(o,r)$ tali che
$(o,r) \notin C_0$ e $(o,r) \in C_D$. Sia data una richiesta che induce $D$ a
elaborare una destinazione per la quale occorre $(o,r)$. L'autorizzazione
invariante rispetto alla linea rileva che $D$ possiede $C_D$, che concede
$(o,r)$; tuttavia, per il teorema di proiezione della
Sezione~\ref{sec:impossibility}, eventi che condividono $(o,r)$ ma differiscono
per linea causale ricevono lo stesso giudizio. La policy non può quindi dipendere
dal fatto che $C_D$ appartenga o meno alla linea della richiesta in esame. Il
sistema ammette pertanto un'esecuzione nella quale la richiesta induce $D$ a
esercitare $(o,r)$ nonostante $(o,r) \notin C_0$, violando la sicurezza rispetto
al delegato confuso: una contraddizione.
\end{proof}

Il teorema non costituisce una critica al possesso: mostra che il possesso
invariante rispetto alla linea è insufficiente per la propagazione multipassaggio
dell'autorità proprio quando si richiedono sia la commistione di autorità sia la
sicurezza rispetto al delegato confuso. Poiché le tre proprietà non possono
valere insieme, un sistema deve rinunciare ad almeno una di esse. Per i sistemi
che richiedono sia una commistione di autorità concretamente utilizzabile, sia
la sicurezza rispetto al delegato confuso, la scelta praticabile è abbandonare
l'invarianza rispetto alla linea, rendendo l'autorizzazione sensibile a essa.
Vietare la delegazione che consente la commistione impedirebbe invece agli
esecutori di detenere autorità propria insieme a quella derivata dalla richiesta:
una soluzione impraticabile nelle installazioni reali, dove un servizio conserva
legittimamente i propri privilegi. Accettare le esecuzioni con delegato confuso
significherebbe rinunciare alla proprietà di sicurezza in esame. La Prova di
Continuità sceglie la prima opzione: non elimina il possesso, ma ripristina la
coordinata temporale mancante, permettendo all'autorizzazione di verificare se
ogni uso di autorità costituisca una continuazione valida del contesto che lo ha
determinato. Il contributo non consiste nel rifiutare il possesso, bensì nel
mostrare che esso è la proiezione spaziale di una relazione di autorità più
ricca; la coordinata mancante è la linea causale temporale.

\begin{table}[ht]
\centering
\small
\begin{tabular}{@{}p{2.9cm}p{5.0cm}p{3.1cm}@{}}
\toprule
Proprietà abbandonata & Conseguenza & Esito \\
\midrule
Delegazione con commistione di autorità & gli esecutori non possono detenere autorità indipendente dalla richiesta & respinta (impraticabile) \\
Sicurezza rispetto al delegato confuso & diventano ammissibili esecuzioni con delegato confuso & respinta (non sicura) \\
Invarianza rispetto alla linea & l'autorizzazione diventa sensibile alla linea & adottata (Prova di Continuità) \\
\bottomrule
\end{tabular}
\caption{Le tre proprietà non possono valere contemporaneamente, perciò un
sistema deve rinunciare ad almeno una di esse. Per i sistemi che richiedono una
commistione di autorità concretamente utilizzabile e la sicurezza rispetto al
delegato confuso, la scelta praticabile è rinunciare all'invarianza rispetto alla
linea: l'opzione adottata dalla Prova di Continuità.}
\label{tab:tradeoff}
\end{table}

\begin{corollary}[La continuità ripristina la sicurezza rispetto al delegato
confuso]
Abbandonando l'invarianza rispetto alla linea si recupera la sicurezza: con la
propagazione non espansiva $C_{i+1} \subseteq C_i$, dunque $C_k \subseteq C_0$,
e un $(o,r) \notin C_0$ non può comparire in alcun passaggio successivo. È
l'impossibilità stabilita nella Sezione~\ref{sec:cdimposs}.
\end{corollary}

\section{Proprietà di sicurezza}
\label{sec:openpassthrough}

\begin{definition}[Autorità limitata dall'origine]
Un modello garantisce un'autorità limitata dall'origine se ogni privilegio
eseguibile al passaggio $n$ era stato concesso in origine al passaggio $0$.
\end{definition}

Affiniamo la Definizione~\ref{def:confused} per distinguere la delegazione
legittima dal delegato confuso. La variabile discriminante non è la fonte
dell'autorità, ma il fatto che la destinazione dell'azione sia designata dalla
richiesta e che il richiedente fosse autorizzato a determinarla.

\begin{definition}[Autorità di causare]
Il contesto iniziale $C_0$ registra l'autorità che il richiedente è autorizzato
a \emph{causare}, non soltanto a esercitare direttamente. Un'azione successiva
appartiene alla linea causale del richiedente se e solo se realizza un'autorità
in $C_0$; altrimenti è un'azione propria del delegato, radicata in un'origine
distinta (il delegato stesso) e non vincolata da $C_0$.
\end{definition}

\begin{definition}[Tramite aperto]
Un delegato è un \emph{tramite aperto} quando, in conseguenza di una richiesta,
applica autorità a una destinazione designata dalla richiesta stessa per
un'azione che il richiedente non era autorizzato a causare:
$(o,r) \notin C_0$, ma $(o,r)$ viene comunque eseguito a causa della richiesta.
Questo è esattamente il delegato confuso.
\end{definition}

Le definizioni separano tre casi. (i) Le attività interne del delegato --- il
compilatore che scrive il proprio file di statistiche --- hanno origine nel
delegato stesso e non sono vincolate da $C_0$. (ii) La delegazione legittima ---
un gateway di pagamento che trasferisce fondi perché l'utente era autorizzato a
causare il pagamento --- realizza un'autorità in $C_0$ ed è una continuazione
valida, sebbene l'utente non detenga direttamente l'autorità bancaria. (iii) Il
delegato confuso --- la scrittura di un file designato dall'utente, che questi
non poteva causare --- è un tramite aperto: l'unico caso vietato dal modello. La
domanda rilevante non è se l'origine potesse compiere direttamente l'operazione
a valle, ma se fosse autorizzata a \emph{causarla} entro questa linea. Quando
servizi successivi usano spazi di operazioni differenti, interviene la
traduzione di policy $\mathcal{T}$ introdotta sopra, la cui sintesi e verifica
sono rinviate a lavori futuri.

\paragraph{Esempio.}
Se $C_0 = \{(\text{convert},r)\}$, allora $(\text{write},r)$ non può mai essere
autorizzato entro la medesima catena di esecuzione:
$(\text{write},r)\notin C_0$ e $C_{i+1}\subseteq C_i$ a ogni transizione,
pertanto $(\text{write},r)\notin C_n$.

\section{Impossibilità del delegato confuso}
\label{sec:cdimposs}

La condizione del delegato confuso richiede un disallineamento di autorità: un
privilegio assente dal contesto di autorità della richiesta viene nondimeno
esercitato in conseguenza della richiesta stessa.

È il secondo risultato d'impossibilità del lavoro. La
Sezione~\ref{sec:impossibility} mostra che il possesso \emph{non può} distinguere
le occorrenze; qui mostriamo il risultato duale: in presenza della continuità,
il delegato confuso \emph{non è uno stato rappresentabile dal modello}. La
condizione non è soddisfacibile in alcuna esecuzione PIC valida: si tratta di
un'impossibilità per costruzione, al livello dei comportamenti validi del
modello.

\begin{theorem}[Il delegato confuso è impossibile nel modello PIC]
In ogni esecuzione PIC valida --- cioè un'esecuzione le cui transizioni
adiacenti soddisfano $PoR$ e la monotonia $C_{i+1}\subseteq C_i$, cosicché
l'autorità è limitata dall'origine ($C_n \subseteq C_0$) --- le condizioni del
delegato confuso non possono essere soddisfatte congiuntamente come
comportamento valido del modello.
\end{theorem}

\begin{proof}
Per $PoR$, i passaggi formano un'unica linea causale radicata in $s_0$ e, per
monotonia, $C_k \subseteq C_0$ per ogni $k$; dunque $(o,r)\notin C_0$ implica
$(o,r)\notin C_k$. L'autorità detenuta dal delegato al di fuori di questa linea
ha un'origine distinta (Sezione~\ref{sec:openpassthrough}) e non viene esercitata
come parte della richiesta; entro la linea della richiesta, il disallineamento
di privilegi necessario al delegato confuso è pertanto insoddisfacibile.
\end{proof}

\paragraph{Interpretazione.}
In PIC l'autorità non è qualcosa che un soggetto semplicemente \emph{possiede};
è una proprietà di continuità della catena di esecuzione. Poiché nessun
passaggio può introdurre o reinterpretare un'autorità assente all'origine,
un'esecuzione con delegato confuso può ancora essere tentata sul piano fisico,
ma non può essere accettata come comportamento valido nel modello PIC.

\section{Discussione}

I sistemi PoP fanno intenzionalmente discendere l'autorizzazione dal controllo
di un artefatto, proprietà spesso utile e necessaria per le decisioni locali di
autorizzazione. In un'esecuzione multipassaggio, tuttavia, il solo controllo
dell'artefatto non stabilisce se un uso successivo sia una continuazione valida
della linea di autorità che lo ha determinato. I sistemi PIC/PoC propagano invece
soltanto sottoinsiemi non espansivi dell'autorità lungo una catena di esecuzione
verificata causalmente.

La Prova di Possesso dimostra il controllo dell'artefatto.
La Prova di Relazione dimostra la relazione fra esecuzioni.
La Prova di Continuità dimostra la continuazione dell'autorità.
La Prova di Possesso stabilisce il controllo su un artefatto di autorità in un
punto dell'esecuzione; da sola, non stabilisce che un uso successivo costituisca
una continuazione valida della linea di autorità che lo ha determinato: è questa
la condizione temporale aggiunta dalla Prova di Continuità. La Prova di Possesso
può fungere da meccanismo per costruire una Prova di Relazione --- il possesso di
un artefatto di continuazione emesso dal predecessore può attestare un singolo
passaggio ---, ma non è l'unico meccanismo possibile.

\paragraph{L'autorità è continuità.}
Nel modello PIC, l'autorità non è una proprietà statica che un'entità si limita
a \emph{possedere}, bensì una proprietà di continuità dell'esecuzione: un
privilegio è autorizzato in un passo se e solo se è una continuazione non
espansiva del contesto che lo ha determinato. In questo senso l'autorità
\emph{è} continua: esiste soltanto come prosecuzione ininterrotta di un'origine,
non come possesso separato dalla propria linea causale.

\paragraph{Spazio delle soluzioni progettuali.}
Il vincolo temporale può essere realizzato in due modi: trasportato
\emph{all'interno} dello stesso artefatto di autorità, come nelle capability
attenuate e vincolate al contesto, oppure mantenuto \emph{accanto} a esso sotto
forma di prova esplicita di continuità. Entrambe le soluzioni possono essere
sensibili alla continuità: per il Teorema~\ref{thm:reintro}, ciò che conta non è
dove risieda il vincolo, ma se la decisione di autorizzazione lo legga. Un
artefatto il cui vincolo venga ignorato --- perché presentato al di fuori
dell'occorrenza che lo ha trasportato --- ricade invece nel possesso invariante
rispetto alla linea. PIC adotta la seconda impostazione, trattando la continuità
come proprietà di prima classe della catena di esecuzione; i due approcci sono
complementari, non alternativi.

L'invariante di sicurezza risultante è semplice: nessun servizio a valle può
compiere, come parte di una catena di esecuzione valida, un'operazione non
autorizzata all'origine. Si tratta di un affinamento del modello, non di una
critica ai meccanismi esistenti: possesso, capability, token al portatore e
vincoli di prova del possesso restano utili strumenti locali di autorizzazione,
mentre PoC rende esplicita l'ulteriore condizione temporale necessaria quando
l'autorità viene propagata lungo una linea causale di esecuzione.

\paragraph{L'autorità come transazione distribuita.}
Una catena multipassaggio di autorità è una transazione distribuita
sull'autorità, anziché sui dati. Un artefatto posseduto consente una decisione
locale di autorizzazione nel punto d'uso --- spesso necessaria per
l'installabilità, le prestazioni e il basso accoppiamento ---, ma revoca, stato
della delegazione, variazioni dell'ambito e linea causale sono fatti globali
riguardanti l'autorità. Autorizzazione locale e coerenza globale dell'autorità
divergono quando tali fatti non sono rappresentati nel punto d'uso. PoC affronta
il problema trasportando esplicitamente la linea causale e imponendo una
propagazione non espansiva dell'autorità lungo la catena.

La continuità non coincide quindi con la propagazione dell'identità. Fornitori
di identità, wallet, credenziali verificabili~\cite{vcdm2022}, flussi sul modello
di OIDC/OIDC4VC~\cite{openidconnect2014,openid4vc2024} e identità dei carichi di
lavoro sul modello di WIMSE/SPIFFE~\cite{wimse2024,spiffe2017} restano importanti
per autenticazione, presentazione delle credenziali, audit, attribuzione di
responsabilità e inizializzazione dell'origine; tuttavia, la ripetizione di un
identificatore nei passaggi successivi non dimostra, di per sé, che l'autorità
esercitata appartenga alla stessa linea causale di esecuzione. Una linea ha
inizio quando un'entità dotata di permessi esprime un intento entro i propri
privilegi, creando un contesto di autorità iniziale $C_0$; i passaggi successivi
non devono necessariamente trasportare o reinterpretare l'identità dell'origine,
ma devono conservare un contesto di autorità non espansivo e causalmente legato
a essa.

\paragraph{Applicazione agli agenti di IA.}
Dal punto di vista della propagazione dell'autorità, gli agenti di IA sono un
caso particolare di esecuzione distribuita autonoma; li menzioniamo
esplicitamente perché in questo contesto la lacuna temporale è particolarmente
marcata. Un agente esegue una catena di passi --- pianificazione, invocazioni di
strumenti, chiamate a valle --- mentre dispone contemporaneamente di più fonti
di autorità: le proprie credenziali di servizio, token utente delegati e ambiti
specifici dei singoli strumenti. La classica designazione del nome di un file
locale riappare come parametro influenzato dall'agente: un identificatore di
tenant, documento o ruolo, un messaggio in coda o la destinazione di una
richiesta. Quando il chiamante influenza tale destinazione e l'agente soddisfa
la richiesta usando autorità proveniente da un contesto diverso, si ottiene il
problema del delegato confuso in forma distribuita. PIC lega ogni invocazione di
strumento e ogni richiesta a valle alla linea causale che l'ha determinata;
pertanto, un'autorità applicata al di fuori dell'esecuzione che ha causato
l'azione non costituisce una continuazione valida. Un agente incaricato di
riassumere un documento contenente una prompt injection indiretta --- per
esempio istruzioni volte a cancellare file o a esfiltrare segreti --- può tentare
tali azioni, ma in PoC esse sono valide soltanto se il contesto di autorità
iniziale contiene i privilegi corrispondenti: se autorizza soltanto
$(\text{summarize},d)$, allora né $(\text{delete},r)$ né
$(\text{exfiltrate},r)$ possono costituire continuazioni valide.

\section*{Lavori correlati}

Il problema del delegato confuso è stato introdotto da
Hardy~\cite{hardy1988confused}. Il controllo degli accessi basato su
capability~\cite{dennis1966programming} e l'analisi del delegato confuso nel
contesto delle capability~\cite{miller2003capability} affrontano la dimensione
spaziale dell'autorità: se un attore detenga autorità su un oggetto. Questo
lavoro isola la dimensione temporale, a essa ortogonale: se l'autorità impiegata
appartenga all'esecuzione che ha causato l'azione. Consideriamo il delegato
confuso non come lo specifico scenario del compilatore di sistema operativo nel
quale fu descritto per la prima volta, ma come il disallineamento generale fra
autorità e causalità esemplificato da quello scenario, e lo riformuliamo per
l'esecuzione discreta in più passi: il contesto dei sistemi distribuiti, delle
catene di servizi e degli agenti di IA, anziché di una singola invocazione
locale. In questa formulazione, il delegato confuso è un'istanza di una classe di
vulnerabilità nella propagazione dell'autorità accomunate dallo stesso
disallineamento; con la Prova di Continuità, tale disallineamento non può
presentarsi come comportamento valido del modello, poiché un privilegio assente
dal contesto di autorità iniziale non può essere esercitato successivamente
nella medesima linea causale. Il lavoro confronta la propagazione dell'autorità
basata sul possesso e quella basata sulla continuità della provenienza
esattamente lungo questo asse.

Il precedente più vicino è il controllo degli accessi basato sulla
storia~\cite{abadi2003history}, nel quale l'autorità disponibile a una
computazione dipende dalla storia di quanto è stato eseguito. Condividiamo il
presupposto che l'autorità sia funzione dell'esecuzione, anziché del possesso
statico. Diversamente da una verifica retrospettiva sui chiamanti precedenti
(come nell'ispezione dello stack), PoC è una disciplina di propagazione in
avanti: ogni passaggio deve essere legato causalmente al predecessore e non deve
espandere l'autorità ricevuta, anche fra passi asincroni e oltre i confini dei
sistemi. Il contributo qui proposto è ortogonale e più specifico: formuliamo la
dipendenza come una \emph{dimensione} $\mathcal{L}$ ortogonale al privilegio
$(o,r)$, dimostriamo che il possesso (ogni policy invariante rispetto a
$\mathcal{L}$) non può distinguere le occorrenze e mostriamo che il delegato
confuso è una collisione nella proiezione che dimentica $\mathcal{L}$. Ne segue
che le mitigazioni basate sul possesso capaci di distinguere tali casi devono
divenire sensibili alla linea causale e, quindi, alla continuità. La letteratura
sulle capability di oggetto~\cite{miller2006robust} analizza il delegato confuso
come separazione fra designazione e autorità, elementi che le capability
riuniscono: questa è la soluzione spaziale. Il nostro risultato aggiunge l'asse
temporale: nell'astrazione del lavoro, una capability trattata soltanto come
artefatto portabile e posseduto, senza vincoli all'occorrenza o alla linea
causale, ricade nella proiezione invariante rispetto alla linea e non può
distinguere due occorrenze che differiscono soltanto per essa.

Il controllo degli accessi basato sulla
provenienza~\cite{park2012provenance} usa la provenienza registrata degli oggetti
come dato in ingresso alle decisioni di autorizzazione; PIC rende invece la
linea causale dell'\emph{esecuzione} stessa il vettore di un contesto di autorità
non espansivo. Il controllo decentralizzato dei flussi di
informazione~\cite{myers1997decentralized} propaga etichette sui dati nel
rispetto di una restrizione monotona; la Prova di Continuità propaga l'autorità
lungo una linea causale di esecuzione secondo la medesima disciplina di
monotonia, ma cambia l'oggetto propagato: l'autorità anziché le etichette dei
flussi di informazione.

\paragraph{Token al portatore, macaroon e vincoli di prova del possesso.}
OAuth 2.0 comprende l'uso di token al portatore come schema comune di
autorizzazione~\cite{hardt2012oauth,jones2012bearer}. Un token al portatore
impiegato senza un vincolo esplicito al contesto o alla linea causale consente
una decisione che legge il token e $(o,r)$, ma non la linea causale del suo uso:
è invariante rispetto a essa e ricade quindi nel primo teorema. I
macaroon~\cite{birgisson2014macaroons} aggiungono \emph{restrizioni} che possono
soltanto attenuarsi e concatenarsi crittograficamente --- esattamente la nostra
attenuazione monotona $C_{i+1}\preceq C_i$ ---, ma un macaroon resta una
credenziale al portatore: in assenza di una restrizione che lo leghi al contesto
della richiesta, può essere presentato fuori dalla linea che lo ha prodotto.
DPoP~\cite{fett2023dpop} e Token Binding~\cite{popov2018token} vincolano un token
a una chiave o a un canale di trasporto; nei termini del
Teorema~\ref{thm:reintro}, tali vincoli fanno sì che la decisione legga un
ulteriore contesto dell'occorrenza e sono quindi localmente sensibili alla
continuità. Ciascun meccanismo riduce l'uso improprio fra contesti nella misura
in cui una restrizione o un vincolo induce la decisione a leggere $\ell$;
nessuno, da solo, stabilisce una Prova di Continuità globale, poiché la decisione
non è ancora tenuta a leggere l'intera linea causale dall'origine lungo una
catena multipassaggio. PIC rende globale ed esplicita tale dimensione.

Il modello PIC è sviluppato indipendentemente in~\cite{gallo2025pic}, dove la
Provenance Identity Continuity viene introdotta come modello generale per i
sistemi di esecuzione distribuiti. Il presente lavoro estende una versione
precedente~\cite{gallo2025apm}, affinandola mediante la distinzione fra relazione
e continuità e l'analisi spaziale e temporale dell'autorità.

\section{Limiti e sviluppi futuri}

Il lavoro presenta un modello lineare minimale della Prova di Continuità. Molte
delle direzioni indicate di seguito non sono difetti del modello, ma proprietà
da sviluppare entro il medesimo principio di continuità; distinguiamo pertanto i
limiti effettivi dell'ambito dalle estensioni naturali. In primo luogo, la revoca
non viene trattata come operazione retroattiva su una linea causale già emessa:
un sistema concreto deve stabilire se essa blocchi soltanto le transizioni
future, invalidi le continuazioni ancora disponibili oppure determini
l'annullamento dell'intera linea. In secondo luogo, la composizione: il modello
è lineare, mentre biforcazioni, ricongiungimenti, fan-out, fan-in ed esecuzioni a
forma di DAG costituiscono una naturale generalizzazione di $\mathcal{L}$ a un
ordine parziale --- un'estensione del principio di continuità a linee composte,
non un suo limite. In terzo luogo, la costruzione di $PoR$ è astratta: la sua
realizzazione concreta è un problema implementativo distinto dal modello di
continuità. L'evidenza di relazione richiesta può essere ricavata da meccanismi
quali ambienti di esecuzione fidati o credenziali verificabili, con proprietà di
sicurezza e prestazioni dipendenti dall'implementazione; il modello è neutrale
rispetto alla scelta.
In quarto luogo, la validazione della linea causale è un problema implementativo,
non un limite del modello: le realizzazioni concrete non devono necessariamente
rivalidare l'intera storia a ogni passaggio e possono usare checkpoint, prove
sintetiche, codifiche compatte o costruzioni basate su accumulatori per mantenere
limitato il costo della validazione.
In quinto luogo, l'origine della linea causale: il modello vincola l'autorità
\emph{entro} una linea, ma non disciplina di per sé quando un esecutore possa
legittimamente aprirne una nuova, radicata nei propri privilegi. Un esecutore
influenzato da un avversario potrebbe classificare erroneamente come autonoma
un'azione causata da una richiesta, compiendola fuori dalla catena della vittima
senza espandere alcun $C_i$. Distinguere le azioni causate da richieste da quelle
che hanno origine nell'esecutore compete alla costruzione di $PoR$ e
all'architettura di applicazione, ed è lasciato a lavori futuri.
Infine, i sistemi eterogenei richiedono traduzioni di policy $\mathcal{T}$ fra
vocabolari locali; ne rinviamo a lavori futuri la sintesi, la verifica e la
governance.

\section{Conclusioni}

Abbiamo introdotto la Prova di Continuità come modello temporale minimale per la
propagazione dell'autorità. La distinzione fondamentale è fra possesso, che
dimostra il controllo di un artefatto di autorità; relazione, che dimostra il
legame causale fra passi di esecuzione adiacenti; e continuità, che dimostra come
l'autorità si sia propagata senza espandersi lungo l'intera linea causale. Il
risultato principale è che il possesso invariante rispetto alla linea non può
distinguere occorrenze che condividono il medesimo privilegio ma discendono da
cause diverse; viceversa, con la Prova di Continuità, un privilegio assente dal
contesto di autorità iniziale non può comparire in seguito nella medesima catena
di esecuzione valida. La condizione del delegato confuso non viene quindi
accettata come comportamento valido del modello: essa è precisamente il
disallineamento fra autorità e causalità escluso dalla continuità. Più in
generale, PoC riformula la propagazione dell'autorità come proprietà temporale di
una linea causale di esecuzione, anziché come mera proprietà statica del possesso
di un artefatto.

\paragraph{Nota sul nome PIC.}
In Provenance Identity Continuity, il termine Identity indica l'identità di
sicurezza o di esecuzione dell'entità dotata di permessi dalla quale l'autorità
viene creata o ancorata all'origine. Non significa che l'identità venga
propagata come primitiva di autorizzazione a ogni passaggio. Il modello sposta
invece l'onere dell'autorizzazione dalla reiterata interpretazione dell'identità
alla verifica della continuità dell'esecuzione: una linea causale che trasporta
un contesto di autorità non espansivo. Provenance indica l'origine e la linea
causale dell'esecuzione; Identity, l'origine autenticata e dotata di permessi;
Continuity, la conservazione senza espansione dell'autorità nei passaggi
successivi. Identità e identificatori restano pertanto essenziali per
autenticazione, presentazione delle credenziali, audit, attribuzione di
responsabilità e inizializzazione dell'origine; PoC riguarda la continuità
dell'autorizzazione che deve valere dopo che tale origine è stata stabilita.

\section*{Ringraziamenti}

Questo lavoro fa seguito a una versione precedente del modello PIC, resa
disponibile su Zenodo nel dicembre 2025~\cite{gallo2025pic}. L'autore ringrazia
Alan H. Karp, la cui esperienza è stata preziosa per mettere alla prova i
concetti di Prova di Continuità e Prova di Relazione. I concetti qui presentati
sono propri dell'autore e il ringraziamento non implica l'approvazione del
lavoro da parte sua. L'autore si è inoltre avvalso di strumenti automatizzati di
assistenza linguistica per l'affinamento editoriale, la formattazione LaTeX e il
perfezionamento della notazione. Tutti i contributi concettuali, i modelli e le
dimostrazioni sono opera dell'autore.

\begin{thebibliography}{99}

\bibitem{hardy1988confused}
N. Hardy.
\newblock The Confused Deputy: Or Why Capabilities Might Have Been Invented.
\newblock \emph{ACM SIGOPS Operating Systems Review}, 22(4):36--38, 1988.

\bibitem{dennis1966programming}
J. B. Dennis and E. C. Van Horn.
\newblock Programming Semantics for Multiprogrammed Computations.
\newblock \emph{Communications of the ACM}, 9(3):143--155, 1966.

\bibitem{miller2003capability}
M. S. Miller, K.-P. Yee, and J. S. Shapiro.
\newblock Capability Myths Demolished.
\newblock Technical Report SRL2003-02, Systems Research Laboratory, Johns
Hopkins University, 2003.

\bibitem{abadi2003history}
M. Abadi and C. Fournet.
\newblock Access Control Based on Execution History.
\newblock In \emph{Proceedings of the Network and Distributed System Security
Symposium (NDSS)}, 2003.

\bibitem{miller2006robust}
M. S. Miller.
\newblock Robust Composition: Towards a Unified Approach to Access Control and
Concurrency Control.
\newblock PhD thesis, Johns Hopkins University, Baltimore, MD, USA, 2006.

\bibitem{hardt2012oauth}
D. Hardt, editor.
\newblock The OAuth 2.0 Authorization Framework.
\newblock RFC 6749, Internet Engineering Task Force, October 2012.

\bibitem{jones2012bearer}
M. Jones and D. Hardt.
\newblock The OAuth 2.0 Authorization Framework: Bearer Token Usage.
\newblock RFC 6750, Internet Engineering Task Force, October 2012.

\bibitem{birgisson2014macaroons}
A. Birgisson, J. G. Politz, \'U. Erlingsson, A. Taly, M. Vrable, and
M. Lentczner.
\newblock Macaroons: Cookies with Contextual Caveats for Decentralized
Authorization in the Cloud.
\newblock In \emph{Proceedings of the Network and Distributed System Security
Symposium (NDSS)}, 2014.

\bibitem{fett2023dpop}
D. Fett, B. Campbell, J. Bradley, T. Lodderstedt, M. B. Jones, and D. Waite.
\newblock OAuth 2.0 Demonstrating Proof of Possession (DPoP).
\newblock RFC 9449, Internet Engineering Task Force, September 2023.

\bibitem{popov2018token}
A. Popov, M. Nystr\"om, D. Balfanz, and J. Hodges.
\newblock The Token Binding Protocol Version 1.0.
\newblock RFC 8471, Internet Engineering Task Force, October 2018.

\bibitem{wimse2024}
IETF WIMSE Working Group.
\newblock Workload Identity in Multi-System Environments (WIMSE).
\newblock Internet Engineering Task Force, 2024.
\url{https://datatracker.ietf.org/wg/wimse/}.

\bibitem{spiffe2017}
SPIFFE Project.
\newblock Secure Production Identity Framework for Everyone (SPIFFE)
Specification.
\newblock Cloud Native Computing Foundation, 2017.
\url{https://github.com/spiffe/spiffe}.

\bibitem{openidconnect2014}
OpenID Foundation.
\newblock OpenID Connect Core 1.0.
\newblock 2014. \url{https://openid.net/specs/openid-connect-core-1_0.html}.

\bibitem{openid4vc2024}
OpenID Foundation.
\newblock OpenID for Verifiable Credentials (OpenID4VC).
\newblock Specification, 2024. \url{https://openid.net/sg/openid4vc/}.

\bibitem{didcore2022}
World Wide Web Consortium (W3C).
\newblock Decentralized Identifiers (DIDs) v1.0.
\newblock W3C Recommendation, 2022. \url{https://www.w3.org/TR/did-core/}.

\bibitem{vcdm2022}
World Wide Web Consortium (W3C).
\newblock Verifiable Credentials Data Model v1.1.
\newblock W3C Recommendation, 2022. \url{https://www.w3.org/TR/vc-data-model/}.

\bibitem{gallo2025pic}
N. Gallo.
\newblock PIC Model --- Provenance Identity Continuity for Distributed Execution
Systems.
\newblock Zenodo, 2025. \url{https://zenodo.org/records/17777421}.

\bibitem{gallo2025apm}
N. Gallo.
\newblock Authority Propagation Models: PoP vs PoC and the Confused Deputy
Problem.
\newblock Zenodo, 2025. \url{https://doi.org/10.5281/zenodo.17833000}.

\bibitem{park2012provenance}
J. Park, D. Nguyen, and R. Sandhu.
\newblock A Provenance-Based Access Control Model.
\newblock In \emph{Proceedings of the International Conference on Privacy,
Security and Trust (PST)}, 2012.

\bibitem{myers1997decentralized}
A. C. Myers and B. Liskov.
\newblock A Decentralized Model for Information Flow Control.
\newblock In \emph{Proceedings of the ACM Symposium on Operating Systems
Principles (SOSP)}, 1997.

\end{thebibliography}

\end{document}
